% Generated draft: requires mathematical and source review.
\documentclass[11pt]{article}
\usepackage[margin=25mm]{geometry}
\usepackage{fontspec}
\setmainfont{Latin Modern Roman}
\usepackage{amsmath,amssymb,mathtools,unicode-math}
\setmathfont{Latin Modern Math}
\usepackage{enumitem,needspace,xurl,hyperref}
\setlength{\parindent}{0pt}
\setlength{\parskip}{5pt}
\setlength{\emergencystretch}{3em}
\newcommand{\N}{\mathbb{N}}
\newcommand{\Z}{\mathbb{Z}}
\newcommand{\Q}{\mathbb{Q}}
\newcommand{\R}{\mathbb{R}}
\newcommand{\C}{\mathbb{C}}
\newcommand{\F}{\mathbb{F}}
\newcommand{\A}{\mathbb{A}}
\newcommand{\PP}{\mathbb{P}}
\newcommand{\E}{\mathbb{E}}
\newcommand{\eps}{\varepsilon}
\newcommand{\dd}{\,\mathrm{d}}
\begin{document}
\section*{UnsolvedMath ID: 9400063}
\textbf{Source problem code:} catalogue entry
\subsection{Definitions and mathematical statement}
In the standard normalization, the Selberg class \(\mathcal S\) consists of Dirichlet series \(F(s)=\sum_{n\ge1}a(n)n^{-s}\), \(a(1)=1\), absolutely convergent for \(\Re s>1\), satisfying:
\begin{enumerate}
\item For some integer \(m\ge0\), \((s-1)^mF(s)\) extends to an entire function of finite order.
\item There are \(Q>0\), \(\lambda_j>0\), \(\Re\mu_j\ge0\), and \(|\omega|=1\), with
\[
\Lambda_F(s)=Q^s\prod_{j=1}^{r}\Gamma(\lambda_js+\mu_j)F(s),
\qquad\Lambda_F(s)=\omega\,\overline{\Lambda_F(1-\bar s)}.
\]
\item For every \(\varepsilon>0\), \(a(n)=O_{F,\varepsilon}(n^\varepsilon)\).
\item In a right half-plane,
\[
\log F(s)=\sum_{n\ge2}b(n)n^{-s},\qquad
b(n)=0\text{ unless }n\text{ is a prime power},\qquad
b(n)=O_F(n^\theta)
\]
for some \(\theta<1/2\), with the logarithm normalized at large real \(s\).
\end{enumerate}
The unit series and products are allowed; primitive elements would be a further restriction.

The question asks whether \(\mathcal S\) equals the intended class \(\mathcal A\) of automorphic L-functions, that is, both whether every such \(F\) has an automorphic realization and whether every admissible automorphic L-function satisfies these axioms. The record does not define \(\mathcal A\): standard \(GL_d\) functions over \(\Q\), other representations/base fields, cuspidal versus isobaric/product realizations, local factors, shifts and normalization all matter.

In particular the sole permitted pole location and the Ramanujan coefficient bound cannot be ignored when interpreting all automorphic L-functions. A nonzero imaginary shift of the zeta function has its pole away from \(1\), illustrating why unrestricted automorphic twists are not automatically compatible with the literal Selberg axioms. Conversely a primitive-only automorphic class would not directly match all products in \(\mathcal S\). These are scope qualifications, not a supplied corrected definition of \(\mathcal A\).

The supplied research summary labels the identification open and also mentions Selberg orthogonality. Orthogonality is a related but separate conjecture and is not silently asserted equivalent to the class-equality question.
\subsection{Short English statement}
Is the Selberg class exactly an appropriately normalized automorphic L-function class? Specify products/primitivity, base field, and pole/twist conventions so that the two classes actually have compatible definitions.
\subsection{Sources}
\begin{itemize}
\item Local UnsolvedMath registry, permanent ID \(9400063\), exact source code \texttt{catalogue entry}, ``Is Selberg class of Dirichlet series equal to class of automorphic L-functions.''
\item Wikipedia, \emph{List of unsolved problems in mathematics}, pinned revision 1366636767, Analytic number theory bullet 63, wikitext line 79; \url{https://en.wikipedia.org/w/index.php?title=List_of_unsolved_problems_in_mathematics&oldid=1366636767#Number_theory}. MediaWiki-source extraction; recorded access 2026-07-29.
\item Supplied OPEN-TRIAGE summary additionally points only to Wikipedia's ``Selberg class'' article. No authoritative primary conjecture statement, definition of the automorphic comparison class, or resolution citation is supplied.
\end{itemize}
\end{document}
